% Manuscript-to-copyeditor.tex - to send to copy editor, taking out the Supp
% Material and various un-needed commented text. 6th January 2020.

% Manuscript.tex - revisions for MEPS. Starting 8th November 2019.
% ssmPart2meps.tex - for submission to MEPS. 18th June 2019.

%%%  FOR TYPESETTERS
%%%  This is my latex file for the manuscript, including references and tables,
%%%  and exluding the Supplementary Material.
%%%  A lot of the commands at the start are just
%%%  settings and preamble, so scroll down to \begin{document} for the start of
%%%  the text. Comments start with %, and I took off the line numbers.  Please
%%%  contact me with any queries, especially if you are going to run Latex
%%%  rather than just use the text from this file, though I expect you may just
%%%  copy the text into another editor and proceed from there).
%%%   Andrew.Edwards@dfo-mpo.gc.ca
%%%
%%%  To make the .pdf file, I run:
%%%  pdflatex Manuscript-to-copyeditor  (maybe twice the first time to get the references working)

\documentclass[12pt, letter]{article}
\textheight 213mm
\topmargin -10mm
\addtolength{\textwidth}{1.0in}
\addtolength{\oddsidemargin}{-0.5in}
\usepackage{epsfig}
\usepackage{rotating}           % For sideways table
\usepackage[pagewise]{lineno}
\usepackage{amsmath}       % for \text for x_min, for \dfrac
\usepackage{cancel}        % for \cancel
\usepackage{natbib}
\usepackage{array}         % for fixed with tables
\usepackage[table]{xcolor} % apparently best to load last. For Appendix tables.
\usepackage{forloop}

\usepackage{mathptmx}      % Bolder Times plus has math fonts (as for hake)

% \usepackage{minitoc}       % See
      % https://tex.stackexchange.com/questions/419249/table-of-contents-only-for-the-appendix
% Make the "Part I" text invisible [ignore main text in TOC]
\renewcommand \thepart{}
\renewcommand \partname{}

% \usepackage{tocloft}       % For ToC
% \addtolength{\cftsubsecnumwidth}{6pt} % For spacing after section numbers in
                                    %  Supp Info TOC.
% \makeatletter \renewcommand{\@dotsep}{4.5} \makeatother % for dots, didn't work.

\bibliographystyle{natbib}


% \linenumbers


% \pagestyle{headings}
\newcommand{\eb}{\begin{eqnarray}}
\newcommand{\ee}{\end{eqnarray}}
\newcommand{\xmin}{x_{\mathrm{min}}}
\newcommand{\xmax}{x_{\mathrm{max}}}
\newcommand{\logten}{\log_{\mathrm{10}}}
\newcommand{\logtwo}{\log_{\mathrm{2}}}
\newcolumntype{L}[1]{>{\raggedright\let\newline\\\arraybackslash\hspace{0pt}}p{#1}}    % From http://tex.stackexchange.com/questions/12703/how-to-create-fixed-width-table-columns-with-text-raggedright-centered-raggedlef   for fixed with columns     though I changed m{#1} to p{#1}

\newcommand{\dx}{\mbox{d}x}
% \newcommand{\standHist}{\ref{fig:dataCode/standHist}}
% \newcommand{\fit}{\ref{fig:dataCode/fitting2a}}
\newcommand{\fitMLEhists}{\ref{fig:binFitting/fitMLEhists}}
\newcommand{\fitMLEconfs}{\ref{fig:binFitting/fitMLEconfs}}
\newcommand{\fitMLEhistsXminSixteen}{\ref{fig:binFitting/xmin16/fitMLEhistsXmin16}}
\newcommand{\fitMLEconfsXminSixteen}{\ref{fig:binFitting/xmin16/fitMLEconfsXmin16}}
\newcommand{\fitMLEhistsCutOffSixteen}{\ref{fig:binFitting/xCutOff16/fitMLEhistsCutOff16}}
\newcommand{\fitMLEconfsCutOffSixteen}{\ref{fig:binFitting/xCutOff16/fitMLEconfsCutOff16}}
\newcommand{\MEE}{erpbb17}     % bib citation for Edwards et al. (2017) MEE paper
\newcommand{\wordcount}{\input{\jobname.wc}}  % To automate wordcount
\newcommand{\epssize}{\epsfxsize=6in}  % To fool latexdiff when adding new
                                % figure (trends100.eps)

% \def\fracd#2{\frac{\displaystyle #1}{\displaystyle #2}}

% Had this for Ecology, probably tweak for JAE:
% \bibpunct{(}{)}{,}{a}{}{,}   % see natbib.pdf, p13

\newcommand\onefig[3]{    % label (old filename and dir) is #1, text is #2,
                          % filename is now #3 (figure1.eps etc for
                          % copyeditor version)
  \begin{figure}[tp]
  \begin{center}
   % \includegraphics[width=6in,height=7in,keepaspectratio=TRUE]{#1.eps} \\  % RH much better control
  \epsfxsize=6in
  \epsfbox{#3.eps}
  \end{center}
  \caption{#2 }
  \label{fig:#1}
  \end{figure}
  % \clearpage
}

\newcommand\onefigshift[3]{    % label (old filename and dir) is #1, text is #2,
                               % filename is #3, as above
  \begin{figure}[tp]
  \begin{center}
   % \includegraphics[width=6in,height=7in,keepaspectratio=TRUE]{#1.eps} \\  % RH much better control
  \epsfxsize=6in
  \epsfbox{#3.eps}
  \end{center}
  \vspace{-10mm}
  \caption{#2}
  \label{fig:#1}
  \end{figure}
  % \clearpage
}

\newcommand\isdfig[2]{    % filename is #1, text is #2
  \begin{figure}[tp]
  \begin{center}
  \vspace{-5mm}
   % \includegraphics[width=6in,height=7in,keepaspectratio=TRUE]{#1.eps} \\  % RH much better control
  % \epsfxsize=6.5in
  \epsfbox{../../fitting-size-spectra/codeBinStructure/nSeaFung/IBTS-ISDfigs/IBTS-ISD#1.eps}
  \end{center}
  \vspace{-5mm}
  \caption{#2}
  \label{fig:ISD-#1}
  \end{figure}
  \clearpage
}


\newcommand\simpfig[2]{    % filename is #1, text is #2. No epsfxsize
  \begin{figure}[tp]
  \begin{center}
   % \includegraphics[width=6in,height=7in,keepaspectratio=TRUE]{#1.eps} \\  % RH much better control
  % \epsfxsize=6in
  \epsfbox{#1.eps}
  \end{center}
  \caption{#2 }
  \label{fig:#1}
  \end{figure}
  \clearpage
}


\newcommand\twofig[4]{   % figure #1 under #2, caption text #3, width #4
  \begin{figure}[tp]     %  label will be #1
  \centering
  \epsfxsize=#4 in
%  \epsfysize=3.5in
  \begin{tabular}{c}
  %	\includegraphics[width=6in,height=3.5in,keepaspectratio=TRUE]{#1.eps} \\  % RH much better control
  %	\includegraphics[width=6in,height=3.5in,keepaspectratio=TRUE]{#2.eps}
  \epsfbox{#1.eps}
  \vspace{-15mm}\\
  \epsfbox{#2.eps}
  \end{tabular}
  \caption{#3}
  \label{fig:#1}
  \end{figure}
  \clearpage
}

\newcommand\threefig[4]{    % figure #1 then #2 then #3,
  \begin{figure}[htp]       %  caption text #4, label will be #1
  \centering
  \begin{tabular}{c}
%	\includegraphics[width=6in,height=2.5in,keepaspectratio=TRUE]{#1.eps} \\  % RH much better control
%	\includegraphics[width=6in,height=2.5in,keepaspectratio=TRUE]{#2.eps} \\
%	\includegraphics[width=6in,height=2.5in,keepaspectratio=TRUE]{#3.eps}
  \vspace{-20mm}
  \epsfbox{#1.eps} \\
  \vspace{-20mm}
  \epsfbox{#2.eps} \\
   \vspace{-20mm}
  \epsfbox{#3.eps}
  \end{tabular}
  \caption{#4}
  \label{fig:#1}
  \end{figure}
}


\renewcommand{\baselinestretch}{1.6}
\renewcommand{\refname}{Literature Cited}
%TC:macroword \onefig 1
%TC:macroword \onefigshift 3
%TC:macroword \simpfig 3
%TC:macroword \simpfigshift 3
%TC:macroword \twofig 3



\begin{document}

%TC:ignore

\noindent {\Large \bf Accounting for the bin structure of data removes
  bias when fitting size spectra}

\vspace{7mm}

\noindent {\large Andrew M.~Edwards$^{1,2,\star}$, James P. W. Robinson$^{2,3}$, Julia
  L. Blanchard$^{4}$, \\Julia K. Baum$^{2}$ and Michael J. Plank$^{5,6}$}

\medskip

{\it
\noindent $^1$Pacific Biological Station, Fisheries and Oceans Canada,
  Nanaimo, British Columbia, V9T 6N7, Canada. %  3190 Hammond Bay Road,
\noindent $^2$Department
  of Biology, University of Victoria, Victoria, British
  Columbia, V8W 2Y2, Canada. % PO Box 1700 STN CSC,
\noindent $^3$Lancaster Environment Centre, Lancaster University, Lancaster,
  LA1 4YW, U.K.
\noindent $^4$Institute for Marine and Antarctic Studies,
  University of Tasmania, Hobart, TAS 7001, Australia. % Private Bag 129,
\noindent $^5$School of Mathematics and Statistics, University of Canterbury,
  Christchurch, New Zealand.
\noindent $^6$Te P\={u}naha Matatini, a New Zealand Centre of
  Research Excellence, University of Auckland, Auckland 1011, New Zealand.}

\medskip

\noindent $^\star$Corresponding author email: Andrew.Edwards@dfo-mpo.gc.ca


\medskip

\medskip

\noindent Running page header: Size spectra for binned data % 3-6 words

\medskip

\section*{Abstract}

Size spectra are recommended tools for detecting the
response of marine communities to fishing or to management measures.
A size spectrum succinctly describes how a
property, such as abundance or biomass, varies with body size in a
community.
Required
data are often collected in binned form, such as numbers of individuals
in 1-cm length bins.
Numerous methods have been employed to fit size spectra, but most give
biased estimates when tested on simulated data, and none account for
the data's bin structure (breakpoints of bins).
Here we use eight methods to fit an annual size-spectrum
exponent, $b$, to an example data set (30~years of the North Sea
International Bottom Trawl Survey).
The methods give conflicting conclusions regarding $b$ declining
(the size spectrum steepening) through time, and so any resulting advice
to ecosystem managers will be highly dependent upon the method used.
Using simulated data we show that ignoring the bin structure gives biased
estimates of $b$,
even for high-resolution data. However, our extended likelihood method, that
explicitly accounts for the bin structure, accurately estimates $b$ and its
confidence intervals, even for coarsely collected data.
We develop a novel visualisation method that accounts for the bin
structure and associated uncertainty, give
recommendations concerning different data types, and
create an R package ({\tt sizeSpectra}) to reproduce all results and encourage use of our methods.
This work is also relevant to wider applications where
a power-law distribution (the underlying distribution for a size spectrum) is
fitted to binned data.

\vspace{10mm}
% 3-8 key words, in order
\noindent {\it Key words:} abundance size spectrum, biomass size spectrum,
individual size distribution, ecosystem-based fisheries management,
ecosystem indicators, truncated Pareto distribution, power-law distribution

% MEE says maximum of 6000-7000 words
\medskip
\noindent Word count (including references, tables and figure legends): 7953 % \wordcount
% this is to automate the word count

% Obtain with:
% texcount -1 -sum -relaxed Manuscript-to-copyeditor.tex > Manuscript-to-copyeditor.wc
% texcount -sum -html Manuscript.tex > Manuscript.html  - looks like it's
%  ignoring bibliography, tried some fixes (even latex to get .bbl, and having
%  .bbl in the file now), but didn't work. Manually counted the difference from
%  submitted version and updated .wc file to 7953.

\medskip

\noindent Revised version to \emph{Marine Ecology Progress Series}

\noindent \today

%TC:endignore

\clearpage

\section{Introduction}

Ecosystem-based fisheries management requires general indicators to accurately
assess ecosystem states across space and time \citep{jenn05, srjfg05}.
Size spectra quantify how a property (such as abundance or biomass) varies with
body size in a community \citep{rg96}, leading to a variety of applications.
They are readily fitted to simple body-size data, and have
therefore been widely applied to aquatic ecosystems
(e.g.~\citealt{dpmg04, dgpr05, bblhw12}).
A new application involves using size-spectra models and simulated carbon fluxes
from an ocean-biogeochemistry model to estimate how global benthic communities
may respond to climate change and ocean acidification \citep{kellyg14, ymabjr17}.
And size
spectra are considered candidate ``ecosystem Essential Ocean
Variables'' for monitoring changes of phytoplankton, zooplankton, benthic
invertebrates, krill and fish in the Southern Ocean \citep{consta16}.


The slope of the size spectrum is a
recommended ecosystem indicator because of its apparent responsiveness to
fishing impacts.
\citet{pj05} tested the response to fishing of eight potential fish community
indicators and found that only the slope of the biomass size spectrum reliably
detected the
effects of spatial management measures (partial closure of an area of
the North Sea). \citet{jd05} similarly recommended the use of size spectra to
provide
surveillance of the status of fish communities. They pointed out difficulties in
justifying reference points (targets to be aimed for and limits to be avoided),
and instead recommended reference directions that monitor trends in
size-spectrum slopes through time. Such trends are commonly investigated in
size-spectrum studies (e.g.~\citealt{blan05, bblhw12}).
Using 188 plausible multispecies size-structured fish community
models, \citet{tllcj15} found that the slope of the abundance size spectrum was
the most responsive indicator of the four indicators tested.
Size-spectra parameters
have been selected by the European Union as one of several indicators used to
evaluate sea-floor integrity \citep{europe10, rice12}, and size-spectrum models
continue to be widely applied
(e.g.~\citealt{tfp14, amln16, ssmarp16, wsj18, zcxxr18, mnwb18, va18, zhou19}).
Such applications require
consistent calculation methods, especially when used for providing practical
advice to fisheries (or ecosystem) managers.
However,
most commonly employed methods for fitting
size-spectra have been shown to give biased estimates and inaccurate confidence
intervals for simulated data
\citep{\MEE}, underscoring the need for close evaluation of methods.


Data used in size-spectra studies are often only collected in binned form,
whereby the
`bin structure' (breakpoints of the bins, \citealt{monn16}) is determined during
data collection. For example, \citet{dgpr05} analysed fish lengths collected in
bins of width 0.5~cm, 1~cm and 5~cm. For benthic invertebrates in the
North Sea,
\citet{mj06} sorted individuals into sixteen $\logtwo$ body-mass bins, resulting
in counts and total biomass for each bin.
Similarly, in underwater visual surveys, fish lengths might be recorded to
the nearest 1~cm \citep{mcco16} or 5~cm \citep{rwnz11},
generating size-frequency data of counts per length bin.
For historical data, all that may be available is
digitisation of points from a
published figure and the bin structure has to be
algebraically determined (e.g.~\citealt{edwa11} in a related context).

Alternatively, data are sometimes binned during analysis. For example,
\citet{dpmg04} and \citet{bblhw12} binned data into 5-cm length bins to
calculate size spectra for fish communities in Fiji and the eastern Bering Sea,
respectively. And some size-spectra calculation methods require binning of
data to a
prescribed resolution \citep{\MEE}.

Here we investigate how the bin structure of data affects the fitting of
size spectra, and whether the bin structure needs to be
explicitly accounted for.
To motivate our work and understand issues that can arise when fitting size
spectra to a complex binned data set,
we first apply eight existing methods to fit annual size spectra to 30 years of data from the
North Sea International Bottom Trawl Survey (IBTS).
The data are the numbers of individuals of each species within each length bin
that were caught per hour of trawling.
We use the IBTS data set as an example data set for several reasons: (i) it is an open-access
data set downloadable from the International Council for the Exploration of the
Sea (ICES) website, (ii) it is an extensive quality-controlled data set
consisting of many species across multiple years,
(iii) we could use the data-processing protocols that were fully
described by \citet{ffrr12} for their calculation of the Large Fish Indicator,
(iv) we could use the species-specific length-weight coefficients collated
by \citet{ffrr12}.

For each year of data we estimate the size-spectrum exponent, $b$. Explicitly
this is the exponent of the individual size distribution (ISD) and
is related to the slope of the size spectrum \citep{weke07, \MEE}. We find
that five methods estimate a significant declining trend in $b$, corresponding
to a steepening size spectrum (as can be expected from fishing pressure,
e.g.~\citealt{rg96} and \citealt{dgpr05}). The remaining three methods
find no significant trend in $b$. Thus, the ecological conclusions differ depending on the
method used.

However, none of the methods properly take into account the bin structure of the
data.
They implicitly assign a single length to all counts within
a length bin, rather than acknowledging that counts can represent a range of
values within the length bin. The binned lengths are subsequently converted to
body masses through species-specific length-weight relationships.
We illustrate the problems that occur, which are compounded due to
most of the methods requiring further binning to estimate $b$.
This motivates investigation of methods to properly account for the bin
structure.

The likelihood method is the only method that can be extended to properly deal
with binned data. We therefore develop and test such an extended
method, which deals with (i) the bin structure of length or body-mass data, and
(ii) situations giving non-integer counts of fish. The latter occurs because the
counts (for the IBTS data)
are standardised to represent the number of fish caught per hour of
trawling, and also arises with other types of data (e.g.~\citealt{rwnz11}).
Using simulated body-mass data we find that applying likelihood
without taking into account the bin structure of the data yields biased
estimates of $b$, even for high-resolution data. However, the extended
likelihood method that takes into account the bin
structure is accurate even
for low-resolution data. Hence it may not be worth the effort and expense of
collecting high-resolution data if the bin structure is not accounted for.

We derive the likelihood
function that fully accounts for species-specific length-weight
relationships. This allows the bin structure for length data to be properly
accounted for when converting lengths to weights.
Applying this function to the IBTS data, we find consistently greater
(less negative) estimates of $b$ than for the original likelihood method.
Finally, we develop a new method for plotting such data and the resulting
ISD that explicitly accounts for the
bin structure and associated uncertainty.

Our results strengthen and extend those from \citet{\MEE}, where we recommended the use of
maximum likelihood estimation over the seven other previous methods based on
simulated (rather than real) data. Here, only the likelihood method can be
properly extended to deal with the uncertainty associated with binned data.

We finish with recommendations concerning which version of the likelihood method
to use for particular types of data. To encourage and enable others to implement
our methods, we provide fully documented and functionalised R code in a new R
package {\tt sizeSpectra}. It includes instructions and vignettes for reproducing
all our results, figures and tables (including those in \citealt{\MEE}), and for
applying our methods to new
data. It is available in the Supplementary Material
or can be installed directly from
\noindent \protect{{\tt https://github.com/andrew-edwards/sizeSpectra}}.

\section{Materials and Methods}

\subsection{Demersal fish data}

For our example data set
we obtained data collected by the North Sea IBTS from the ICES online Database
of Trawl Surveys ({\tt
  http://www.ices.dk/marine-data/data-portals/Pages/DATRAS.aspx}).
We downloaded the number of
organisms caught per hour for each species and
length class for all surveys in Areas~1 to~7 carried out in Quarter~1
from 1986-2015 (see Supplementary Material for full details).
All surveys followed standardized
gear- and data-collection protocols \citep{ices15}.
We followed
\citet{ffrr12} in removing all non-demersal fish.
During surveys, lengths were recorded as rounded
down to the nearest 1~cm [or 0.5~cm for Atlantic Herring (\emph{Clupea
  harengus}) and European Sprat (\emph{Sprattus sprattus})].
We averaged counts of the same year/species/length-class combination
across the seven areas.
For simplicity this ignores potential differences in
sampling effort due to variations in haul duration, trawl speed and net area,
and fits a common size spectrum for the whole region.

For each species, $s$, the length-weight relationship is
\eb
w = \alpha_s l^{\beta_s},
\label{lw}
\ee
where $w$ is the estimated body mass (g) of an individual, $l$ is the known
body length (cm), and $\alpha_s$ and $\beta_s$ are species-specific parameters
published by \citet{ffrr12}. Following \citet{ffrr12}, for $l$ we initially use
the `Length class' value (see
Table~\ref{tab:egData} for examples), which represents the minimum of possible
lengths for a given fish (i.e.~a `Length class' of 35~cm represents fish in the
range 35-36~cm). This is consistent with, for example, \citet{blan05} and
\citet{dgpr05}, the latter also doing this for 5-cm and 10-cm length
bins. For example, a 45-cm Smallspotted Catshark is assigned a calculated
body mass of $\alpha_s l^ {\beta_s} = 0.0031 \times 45^{3.0290} = 315.46$~g
(first row of Table~\ref{tab:egData}). In 1986 there were 0.00714~h$^{-1}$ of such fish caught,
resulting in a total biomass per hour of trawling of
2.25~g~h$^{-1}$ (Table~\ref{tab:egData}).

Following calculations of size spectra for earlier IBTS data \citep{pj05} and
Celtic Sea groundfish surveys \citep{blan05}, we remove all resulting body
masses estimated to be $<4$~g because small fish will not be effectively sampled by
the gear. The resulting data set (Table~\ref{tab:egData}) contains the required
information to estimate the size-spectrum exponent of the community for each
year. It consists of a dataframe with 42,298 rows, where each row is a unique
combination of year, species and length class. Each year has between 57 and 81
unique species.

\subsection{Individual size distribution}

We designate the random variable $X$ to represent the mass (g) of an individual fish
or other organism. Then, considering $X$ to come from a bounded power-law (PLB)
distribution, the probability density function for $X$ is \eb f(x) = C
x^b,~~~~\xmin \leq x \leq \xmax
\label{PLB}
\ee where \eb C = \left\{ \begin{array}{ll} \dfrac{b + 1}{\xmax^{~b + 1} -
    \xmin^{~b + 1}}, \vspace{2mm} & b \neq -1, \\ \dfrac{1}{\log \xmax - \log
    \xmin}, & b = -1,
  \end{array}
  \right.
\label{PLB.C}
\ee $x$ represents possible values of $X$, $\log$ is the natural logarithm, $b$
is the size-spectrum exponent, and $\xmin$ and $\xmax$ are the minimum and
maximum possible values of the data with $0 < \xmin < \xmax$ \citep{\MEE}. The
exponent $b$ is expected to be negative, and is the quantity of interest when
analysing size spectra. Equation~(\ref{PLB}) is the ISD,
and the resulting biomass size spectrum (biomass density function) for a
community of $n$ individuals is
\eb
B(x) = n C x^{b+1},~~~~\xmin \leq x \leq \xmax.
\label{biomass}
\ee
For length size spectra $x$ would represent lengths, but (\ref{biomass}) cannot
be directly used \citep{\MEE}.

\subsection{Estimating $b$ for each year of the demersal fish data}

We apply eight methods for fitting size spectra that have been
previously used in publications and in the
R package {\tt mizer} \citep{sba14}. The methods were
documented by \citet{\MEE}. Here, the Supplementary Material explains how
the methods need to be adapted to deal with the fact that the data are not
simply measurements of each individual fish; full results are also given.
Using each method we calculate
the size-spectrum exponent $b$ separately for each year of data.
This gives us a time series of estimated annual values of $b$
with confidence intervals, one time series for each method. We then fit a
weighted linear regression (that uses the confidence intervals) to each
time series to see whether there is a
significant change ($p < 0.05$) in $b$ through time. For clarity we now
present the results which then motivate development of further methods.

\section{Results}

The linear regression results for each method (Figure~\ref{fig:nSeaFung/nSeaFungTrends}) show that
(a) the absolute estimates of
$b$ are highly dependent upon the method used, (b) the confidence intervals
can vary in size through time and are far broader for some methods than
for others,
and (c) there is no significant change in $b$ through time when
using three of the methods
yet a significant negative trend (steepening of the size spectrum)
when using the remaining five methods.
This latter point shows that
methodological differences can lead to
differing ecological conclusions.

\subsection{Dealing with the bin structure of the data}

However, the aforementioned analyses (Figure~\ref{fig:nSeaFung/nSeaFungTrends})
implicitly ignore the fact that using the
minimum of a length
bin, and converting that to a single body mass, does not properly account for
the bin structure of the data (because the assigned length can actually correspond to any true
length within a bin). We demonstrate this issue using length-weight parameters
for relationship (\ref{lw}) for two example species considered by \citet{blan05},
namely $\alpha_s= $0.0255 and $\beta_s=$2.7643 for Lemon Sole ({\it Microstomus
  kitt}) and $\alpha_s=$0.001 and $\beta_s=$3.4362 for Common Ling ({\it Molva
  molva}).


Figure \ref{fig:binFigs/binLW} shows how six 5-cm length bins get converted into
six body-mass bins via (\ref{lw}).
For clarity we use length bins of width 5~cm and use the midpoint of a
length bin to calculate the resulting body
mass (note that the IBTS data mostly uses 1-cm bins and
was analysed above using the minima of length bins rather than midpoints).
The resulting body-mass bins in Figure \ref{fig:binFigs/binLW} are not of equal
width because of the nonlinear nature of the length-weight relationships. And
the body-mass bins differ between species because of the species-specific
length-weight parameters. Such information regarding the body-mass bins was not
explicitly accounted for in the above analyses of the IBTS data; all methods
used a body mass obtained from applying (\ref{lw}) to the minimum of a length
bin.

Figure \ref{fig:binFigs/binLWlog} shows the consequence of using the resulting
body-mass values, converted from the midpoints of the length bins, when using
methods to fit size spectra that require further binning of the data. For
example, the LBbiom
(logarithmic binning to fit biomass size spectrum)
and LBNbiom (logarithmic binning with normalisation to fit biomass size
spectrum)
methods were used by \citet{jd05} and
\citet{pj05}, respectively, to calculate trends in size-spectra slopes for
earlier IBTS data. Explicitly shown is the assignment of the body-mass values
into bins with bin breaks of 4, 8, 16, 32, ..., as would be used in the LBbiom
and LBNbiom methods.

For example, consider blue length-bin number~5 that is highlighted in
Figure~\ref{fig:binFigs/binLW} for Lemon Sole and covers the range
30-35~cm. Say there are 20 individual Lemon Sole that have lengths in this
range. All 20 individuals would be represented by the length equal to the
midpoint of the length bin, namely 32.5~cm, which gets converted to a body-mass
value of 385~g (even though the 20 body masses really lie in the range
309-473~g). This value of 385~g lies within the $\logtwo$ bin that spans
256-512~g (see Bin~5 in Figure~\ref{fig:binFigs/binLWlog}). Therefore, all 20
individuals get assigned a body mass of 384~g, which is the midpoint of the
$\logtwo$ bin. So, in the LBbiom and LBNbiom methods, all 20 individuals with
body masses in the range 309-473~g are assigned body masses of 384~g. The
uncertainty in the original lengths
(which is often an unavoidable consequence of measurement methods)
is not explicitly accounted for.

Furthermore, the length range 20-25~cm (Bin~3 in Figures~\ref{fig:binFigs/binLW}
and~\ref{fig:binFigs/binLWlog}) represents fish in the range 101-187~g.  The
midpoint of 22.5~cm gets converted to a body-mass value of 139~g, which lies
within the $\logtwo$ bin that spans 128-256~g
(Figure~\ref{fig:binFigs/binLWlog}). All individuals then get assigned a body
mass of 192~g (the midpoint of the $\logtwo$ bin), but this is \emph{greater}
than the body mass of any of the individuals (101-187~g).  A similar undesirable
situation occurs for Bin~6 (Figure~\ref{fig:binFigs/binLWlog}).  None of the
remaining bins yield counts in $\logtwo$ bins that accurately represent the
original length bins (since there is no reason that they should).


Similar situations also occur for the remaining four methods that involve
further binning of data.
The LCD (logarithmic plotting of one minus the cumulative distribution) and
MLE (maximum likelihood estimate)
methods do not bin data and so do not
suffer from the extra binning described in Figure~\ref{fig:binFigs/binLWlog},
but do suffer from assigning the same body mass to all individuals in a
length bin.
It is not clear how the LCD method could be extended to account for the bin
structure. The
MLE method is the only method to accurately estimate $b$ and its confidence
intervals for unbinned data \citep{\MEE} --
this motivates our MLEbin
method, which adapts the MLE method to explicitly account for the uncertainty
inherent in binned data.



\subsection{Likelihood methods for binned data}

\subsubsection{MLEmid method}

The simplest way to deal with binned data in a likelihood framework is to just
use the midpoints of the body-mass bins, which we call the MLEmid (maximum
likelihood estimate based on mid-points) method. If there is, say, a count of 10
individuals assigned to the bin that spans the body-mass range 4-8~g, then all
10 individuals are assumed to have a body mass of 6~g (the midpoint, though
recall that for the IBTS data we followed \citealt{ffrr12} in using the minimum value for all eight methods,
including MLE). In
the Supplementary Material we derive the likelihood function for count data in
general and explain how this is used for the MLEmid method.

\subsubsection{MLEbin method}

The MLEbin (maximum likelihood estimate based on binning) method calculates the
MLE for $b$ by explicitly accounting for the fact that each count represent
values within its respective size bin. So for 10 individuals in the body-mass
bin 4-8~g, the MLEbin method uses a likelihood function that explicitly accounts
for the uncertainty in the 10 individual body masses. The body masses can be
anywhere in the range 4-8~g, rather than assuming that they are all 6~g as for
the MLEmid method. The likelihood function for the MLEbin method was derived as
equations (A.70) and (A.75) in \cite{\MEE}.

\subsubsection{Testing the MLEmid and MLEbin methods using simulated data}

We now simulate data to test the two methods. Following \citet{\MEE} we generate
10,000 simulated data sets, each containing 1,000 independent random numbers
(body masses)
drawn from the bounded power-law distribution (\ref{PLB}) with $b=-2$, $\xmin=1~$g
and $\xmax=1,000~$g. We bin each simulated data set to
represent the effect of collecting data in a binned form.
We then test how well
the MLEmid and MLEbin methods estimate $b=-2$ only using the information from
the binned data set (i.e. the bin definitions and the count in each bin), by
maximising the appropriate likelihood function and using the profile
likelihood-ratio test \citep{hm97} to calculate confidence intervals.

We test four binning types. Type `Linear 1' uses bins that all
have a constant bin width of 1~g and that span the data. The first bin starts at
the integer below the lowest data point.
Similarly, `Linear 5' and `Linear 10' use bins of constant widths 5~g
and 10~g, respectively. The binning type `2k' refers to bin widths that double in
size, span the data, and have bin breaks that are integer powers of 2 (i.e. are
at $2^k$~g where $k$ is an integer, giving bin breaks at 1, 2, 4, 8, 16, ...).

Figure~\fitMLEhists~shows that the MLEbin method clearly performs better than
the MLEmid method for all binning types (the histograms are centered around the
true value of $b=-2$); summary statistics are given in the Supplementary
Material. The MLEmid method always overestimates $b$, whereas the
MLEbin gives roughly the same results regardless of the binning (with mean and median
estimates of $b$ of -1.99 or -2.00). Thus the MLEbin method removes the bias
associated with the MLEmid method (note that we cannot formally prove this
is always true as there is no analytical solution for the MLE of $b$, but it
also holds for all simulation results in the Supplementary Material).

Somewhat surprisingly, the MLEmid method is biased even for the `Linear 1'
binning type. So, binning the data using bin breaks of 1, 2, 3, ..., 999, 1000~g
and then using the midpoint of a bin to represent all counts in that bin in a
likelihood context, gives a biased estimate of $b$ (99\% of the simulations
over-estimated $b$). This means that for, say, body-mass data that span the
range 1~g to 1~kg and are measured to a resolution of 1~g, ignoring the fact
that the binned data represent values in a range (i.e. using the MLEmid method)
will lead to biased values of the exponent $b$. The reason for this is the very
skewed nature of the power-law distribution. For example,
for the simulated data set
shown in Figure~S.41 of the
Supplementary Material, the
bin covering the range 1-2~g has a count of 528, but 348 (66\%) of these values
are actually $<1.5$~g; the MLEmid method assumes that all values equal
1.5~g. Whereas the MLEbin method, by definition, accounts for the fact that the
values within a bin are assumed to exhibit a power-law distribution.


For the `Linear 10' binning type, the MLEmid method performs particularly poorly
(Figure~\fitMLEhists(e)). But, perhaps unexpectedly, the corresponding MLEbin method in
Figure~\fitMLEhists(f) leads to accurate estimation of $b$, and furthermore has
better accuracy than the MLEmid method with the finer `Linear 1' binning type in
Figure~\fitMLEhists(a).
Therefore, the MLEbin method outperforms the MLEmid method, even when it
is using data collected at a coarser resolution (10~g vs 1~g) than than the MLEmid method.
This shows that expending additional time and money to collect
high-resolution data
is not
worthwhile if the bin structure is unaccounted for -- it is better to collect
low-resolution data and use the MLEbin method to account for the bin structure.

The `2k' binning type performs slightly better than `Linear 10' when using the
MLEbin method (Figures~\fitMLEhists(f) and~(h)) because it has higher resolution
for the data-rich lower bins (bin breaks are 1, 2, 4, 8, ..., compared to 1, 11,
21, 31, ..., for `Linear 10').

The results regarding confidence intervals (Figure~\fitMLEconfs) give a similar
conclusion. For the MLEmid method and `Linear 1' binning type, only 46\% of the
95\% confidence intervals contain the true value of $b$, less than half of the
desired observed coverage of 95\%. The results are worse as the bin width
increases (`Linear 5' and `Linear 10') and for bin widths that double
(`2k'). However, the MLEbin gives reliable confidence intervals (observed
coverage of 94\% or 95\%) regardless of binning type.

Thus, these results demonstrate that for binned data the MLEmid method is
inaccurate and the MLEbin method should be used, even for seemingly narrow bins.

\subsubsection{The MLEbins method  --  accounting for species-specific body-mass bins}

The IBTS data has species-specific length-weight
relationships, such that the
body-mass bins are different for each species, as occurs for many data sets;
e.g.~\citet{mnwb18} had lengths of 1-cm resolution that were
converted into body masses (and then used for the LBNbiom method).
To account for this we extend the MLEbin
method to the MLEbins method, using a multinomial
log-likelihood function \citep{lawl03}; see Supplementary Material.

\subsection{MLEbins method applied to IBTS data}

To apply the MLEbins method to the IBTS data we first need to calculate
species-specific body-mass bins that result from the species-specific
length-weight coefficients (as for Figure~\ref{fig:binFigs/binLW}).
Figure~\ref{fig:nSeaFung/nSeaFungBins1} shows how the 1-cm (or 0.5-cm) length
bins for 45 of the species get converted to body-mass bins. The conversions are
different for each species because of the different values of the length-weight
coefficients. Thus, the vertical bars are analogous to those in
Figure~\ref{fig:binFigs/binLW}. Figure~\ref{fig:nSeaFung/nSeaFungBins1} shows
that even though the length bins are only 1~cm wide, the resulting individual
body-mass bins can be quite large; the widest bin across all species is 832~g for
Atlantic Cod (\emph{Gadus morhua}) [see Supplementary Material].

In Figure~\ref{fig:ISD-1999.main} we introduce a plotting technique
that visualises the binning structure of the data and the fit of the ISD
(from using the MLEbins method).
The fitted PLB model shows
a reasonable fit to the 1999 data (the red curve passing through the grey
rectangles) for body masses $<100$~g; for larger body masses the PLB model
generally over-estimates the numbers of individuals. This pattern is seen for
many of the years (see Supplementary Material for figures)
and may potentially be a consequence of fishing pressure.


Using the species-specific MLEbins method applied to each year of the IBTS data, we obtain
consistently higher estimates of $b$ than for the original MLE method
(Figure~\ref{fig:nSeaFung/nSeaFungCompareTrends}).
The differing estimates of $b$ show that the likelihood method did need to be
properly adapted (the MLEbins method) to deal with the binned nature of the
data.
In this case, both methods conclude that there is no significant temporal trend
in $b$ through time (see Supplementary Material). Nevertheless, the results show
the potential for different conclusions about community structure to be reached
if the bin structure of data is not properly considered.

\subsection{Recommendations}

In Table~\ref{tab:recommend} we recommend fitting methods based on the type
and resolution of data available.
If data are available in unbinned form (e.g.~\citealt{tdss15})
then the MLE method can be used, but this still depends on the accuracy of the raw
data. Given that unbinned data are
essentially (to some resolution) binned data,
given the results in Figures~\fitMLEhists\ and~\fitMLEconfs\ we recommend
testing the MLE versus MLEbin method for the particular bin structure and range
of a data set.
Or one could just use the MLEbin method, with the bin structure
defined by the accuracy of the measurements.
Our methods may need to be tailored for explicit data sets depending on the
survey design -- in Table~\ref{tab:recommend} we give an example concerning
monitoring of Pacific reefs.

\section{Discussion}

We have analysed a rich 30-year data set from the North Sea, and found that
ecological conclusions regarding apparent trends in size-spectra exponent can depend on
the method used to estimate the exponent. Furthermore, the likelihood method
needs to properly account for the binned nature of data (giving the MLEbin
method). Using simulated data, even for
seemingly high-resolution data the likelihood method that does not account for
the binning (the MLEmid method) is biased, whereas MLEbin is not.
For data, such as the IBTS data, where lengths have to be converted to body
masses using species-specific length-weight coefficients, we have derived
the MLEbins method to estimate $b$ and developed a technique for plotting the
data and the resulting fit.

Following \citet{\MEE} we have used a bounded power-law distribution for the ISD
because power laws are commonly used models for size spectra \citep{pd78, bd92}
and alternative models of size-structured predator-prey dynamics predict
that the aggregate community ISD may indeed be
close to a power law \citep{ab06}.
Future work could focus on relaxing the
assumption of a power-law distribution
if other distributions are deemed feasible \citep{\MEE}.
Using a bounded (rather than unbounded) distribution avoids extrapolation
beyond the range of the data -- this regrettably prevents estimation
of the population density of Loch Ness monsters (as per \citealt{sk72}).
Since power-law distributions have many applications in ecology \citep{weg08},
our simulation results suggest that bin structure may need to be accounted for
in other contexts, and that our novel plotting method
(Figure~\ref{fig:ISD-1999.main}) would be similarly applicable.

Ideally, our methods could be extended to
incorporate estimated uncertainties of the species-specific length-weight
coefficients (see \citealt{froese06}). This would provide further
uncertainty regarding the body mass associated with each length bin
(see \citealt{bbmr19} for an explicit example from raw data).
However, practical applications could be limited since published tabulations of
length-weight relationships do not generally include confidence intervals for the parameters
(e.g., \citealt{lyt02, robins10, ffrr12}).


Following \citet{pj05} and \citet{blan05} we used a (minimum) cutoff of 4~g to
analyse the IBTS data. This was necessarily done by converting the minima of the
length bins into species-specific body-mass bins, and then only using the
resulting bins with minima $\geq$4~g.  For a resulting body-mass bin of
3-5~g, for example, this approach removes some individuals $\geq$4~g, even though we wish to
impose a minimum of 4~g. This may be unavoidable, although note that the
likelihood function for the MLEbins method properly assumes no data (rather than no
individuals) $<$5~g for such a species with lowest bin 3-5~g [$w_{s,1} = 5$~g
for this species in (S.18) in the Supplementary Material].
Also note that we have estimated $\xmax$ separately for each year although it
could be set to a fixed maximum value across all years.

There are techniques available for finding the value of $\xmin$ that gives the
best fit of an unbounded power law \citep{csn09, gilles15}.
But these are focussed on finding power-law \emph{tails} (given the optimum
$\xmin$), rather than finding a distribution that fits the full size range of
the sampled community, which is what is desired when fitting size spectra.
Such techniques would have to be adapted to use our binned likelihood methods (and to
use the bounded rather than unbounded power-law distribution).
Rather, we recommend that the range
of data be justified from biological and sampling perspectives, rather than adjusting
the range to most closely resemble a power-law distribution. This is particularly
so when comparing fits across multiple years or locations because $\xmin$ should be
consistent; e.g.~\citet{robins17} fixed $\xmin = 20$~g for all data sets from 38
Pacific islands.
We have re-run our analyses of the IBTS data with a
cutoff value of 100~g (much higher than the original 4~g). The
exponent $b$ is much lower for all years, and increases through time rather than not
changing (see Supplementary Material). This emphasises the need for careful
biological consideration of the cutoff value in such applications.

We have used the IBTS data to motivate development of widely applicable
methods. We caution against using our current results for management purposes
concerning the North Sea, because we have not fully explored assumptions (such
as those just discussed) and have used all the sampled species rather than only
a well-sampled subset (as recommended by \citealt{jd05}). Also, the power
analyses conducted by \citet{jd05} on earlier IBTS data could be repeated, in
particular to determine the lower cutoff value that would provide the greatest
power to detect trends in the size-spectrum exponent. Note that \citet{jd05}
were restricted to such cutoffs being powers of two because of the use of the
LBbiom method, but this would not be necessary with the MLEbins method.


Note that the confidence intervals calculated for the IBTS data
from the MLEbins method
are not strictly true confidence intervals, because the commonly-used unit of `numbers
of individuals caught per hour of trawling' is arbitrary. If numbers were
instead expressed as `per day of trawling' then all counts would be divided by
24, and
the resulting confidence intervals would be larger than those in
Figure~\ref{fig:nSeaFung/nSeaFungCompareTrends}.
True confidence intervals could be obtained by using the absolute numbers of
individual fish caught summed across all areas (rather than numbers per hour of
trawling averaged across areas); this should be investigated if conducting
further analyses of such types of data.
This would change the width of the confidence intervals, but not the MLE for
$b$, which is unaffected by any change in units.

All results, figures and tables presented here are reproducible in our R package
{\tt sizeSpectra}. The code is properly documented and functionalised. The included
vignettes will enable users to fit size spectra to their own data using our recommended
and tested likelihood methods that properly account for the bin structure of the data.

%TC:ignore
\section*{Supplementary Material}

The following supplementary material is included:
(A)~document containing further methodological details, including derivations
of appropriate likelihood functions, further results for the IBTS and
the simulated data, and a summary of the {\tt sizeSpectra} R package;
(B)~code repository (our new {\tt sizeSpectra} R package) such that all results can be reproduced and users
can apply the methods to their own data.

\section*{Acknowledgements}

We thank Carrie Holt, Brooke Davis, Jaclyn Cleary and Jennifer Boldt for useful
discussions, and Sean Anderson for tutorials on creating R packages. We thank
the two anonymous reviewers for their thoughtful and positive comments that have
helped improve this work.
% TC:endignore


%TC:incbib

\begin{thebibliography}{}

\bibitem[\'Alvarez {\em et~al.}(2016)]{amln16}
\'Alvarez, E., Mor\'an, X. A.~G., L\'opez-Urrutia, A., and Nogueira, E. (2016).
\newblock {Size-dependent photoacclimation of the phytoplankton community in
  temperate shelf waters (southern Bay of Biscay)}.
\newblock {\em Mar.~Ecol.~Prog.~Ser.}, {\bf 543}, 73--87.

\bibitem[Andersen and Beyer(2006)]{ab06}
Andersen, K.~H. and Beyer, J.~E. (2006).
\newblock Asymptotic size determines species abundance in the marine size
  spectrum.
\newblock {\em Am.~Nat.}, {\bf 168}, 54--61.

\bibitem[Benoist {\em et~al.}(2019)]{bbmr19}
Benoist, N. M.~A., Bett, B.~J., Morris, K.~J., and Ruhl, H.~A. (2019).
\newblock A generalised volumetric method to estimate the biomass of
  photographically surveyed benthic megafauna.
\newblock {\em Prog.~Oceanogr.}, {\bf 178}, 102188.

\bibitem[Blanchard {\em et~al.}(2005)]{blan05}
Blanchard, J.~L., Dulvy, N.~K., Jennings, S., Ellis, J.~R., Pinnegar, J.~K.,
  Tidd, A., and Kell, L.~T. (2005).
\newblock Do climate and fishing influence size-based indicators of {C}eltic
  {S}ea fish community structure?
\newblock {\em ICES J.~Mar.~Sci.}, {\bf 62}, 405--411.

\bibitem[Boldt {\em et~al.}(2012)]{bblhw12}
Boldt, J.~L., Bartkiw, S.~C., Livingston, P.~A., Hoff, G.~R., and Walters,
  G.~E. (2012).
\newblock Investigation of fishing and climate effects on the community size
  spectra of eastern {B}ering {S}ea fish.
\newblock {\em Trans.~Am.~Fish.~Soc.}, {\bf 141}, 327--342.

\bibitem[Boudreau and Dickie(1992)]{bd92}
Boudreau, P.~R. and Dickie, L.~M. (1992).
\newblock Biomass spectra of aquatic ecosystems in relation to fisheries yield.
\newblock {\em Can. J. Fish. Aquat. Sci.}, {\bf 49}, 1528--1538.

\bibitem[Clauset {\em et~al.}(2009)]{csn09}
Clauset, A., Shalizi, C.~R., and Newman, M. E.~J. (2009).
\newblock Power-law distributions in empirical data.
\newblock {\em SIAM Rev.}, {\bf 51}, 661--703.

\bibitem[Constable {\em et~al.}(2016)]{consta16}
Constable, A.~J., Costa, D.~P., Schofield, O., Newman, L., {Urban Jr.}, E.~R.,
  Fulton, E.~A., Melbourne-Thomas, J., Ballerini, T., Boyd, P.~W., Brandt, A.,
  de~la Mare, W.~K., Edwards, M., El\'eaume, M., Emmerson, L., Fennel, K.,
  Fielding, S., Griffiths, H., Gutt, J., Hindell, M.~A., Hofmann, E.~E.,
  Jennings, S., La, H.~S., Mc{C}urdy, A., Mitchell, B.~G., Moltmann, T.,
  Muelbert, M., Murphy, E., Press, A.~J., Raymond, B., Reid, K., Reiss, C.,
  Rice, J., Salter, I., Smith, D.~C., Song, S., Southwell, C., Swadling, K.~M.,
  {Van de Putte}, A., and Willis, Z. (2016).
\newblock {Developing priority variables (``ecosystem Essential Ocean
  Variables'' -- eEOVs) for observing dynamics and change in Southern Ocean
  ecosystems}.
\newblock {\em J.~Mar.~Syst.}, {\bf 161}, 26--41.

\bibitem[Daan {\em et~al.}(2005)]{dgpr05}
Daan, N., Gislason, H., Pope, J.~G., and Rice, J.~C. (2005).
\newblock Changes in the {N}orth {S}ea fish community: evidence of indirect
  effects of fishing?
\newblock {\em ICES J.~Mar.~Sci.}, {\bf 62}, 177--188.

\bibitem[Dulvy {\em et~al.}(2004)]{dpmg04}
Dulvy, N.~K., Polunin, N. V.~C., Mill, A.~C., and Graham, N. A.~J. (2004).
\newblock Size structural change in lightly exploited coral reef fish
  communities: evidence for weak indirect effects.
\newblock {\em Can. J. Fish. Aquat. Sci.}, {\bf 61}, 466--475.

\bibitem[Edwards(2011)]{edwa11}
Edwards, A.~M. (2011).
\newblock Overturning conclusions of {{L}}\'evy flight movement patterns by
  fishing boats and foraging animals.
\newblock {\em Ecology}, {\bf 92}, 1247--1257.

\bibitem[Edwards {\em et~al.}(2017)]{erpbb17}
Edwards, A.~M., Robinson, J. P.~W., Plank, M.~J., Baum, J.~K., and Blanchard,
  J.~L. (2017).
\newblock Testing and recommending methods for fitting size spectra to data.
\newblock {\em Meth.~Ecol.~Evol.}, {\bf 8}, 57--67.

\bibitem[{European~Commission}(2010)]{europe10}
{European~Commission} (2010).
\newblock {Commission decision of 1 September 2010 on criteria and
  methodological standards on good environmental status of marine waters}.
\newblock {\em Off. J. Eur. Union}, {\bf L232}, 14--24.

\bibitem[Froese(2006)]{froese06}
Froese, R. (2006).
\newblock Cube law, condition factor and weight-length relationships: history,
  meta-analysis and recommendations.
\newblock {\em J. Appl. Ichthyol.}, {\bf 22}, 241--253.

\bibitem[Fung {\em et~al.}(2012)]{ffrr12}
Fung, T., Farnsworth, K.~D., Reid, D.~G., and Rossberg, A.~G. (2012).
\newblock Recent data suggest no further recovery in {N}orth {S}ea {L}arge
  {F}ish {I}ndicator.
\newblock {\em ICES J.~Mar.~Sci.}, {\bf 69}, 235--239.

\bibitem[Gillespie(2015)]{gilles15}
Gillespie, C.~S. (2015).
\newblock {Fitting heavy tailed distributions: The poweRlaw package}.
\newblock {\em J.~Stat.~Softw.}, {\bf 64}, 1--16.

\bibitem[Heenan {\em et~al.}(2016)]{heen16}
Heenan, A., Gorospe, K., Williams, I., Levine, A., Maurin, P., Nadon, M.,
  Oliver, T., Rooney, J., Timmers, M., Wongbusarakum, S., and Brainard, R.
  (2016).
\newblock Ecosystem monitoring for ecosystem-based management: using a
  polycentric approach to balance information trade-offs.
\newblock {\em J.~Appl.~Ecol.}, {\bf 53}, 699--704.

\bibitem[Hilborn and Mangel(1997)]{hm97}
Hilborn, R. and Mangel, M. (1997).
\newblock {\em The Ecological Detective: Confronting Models with Data}.
\newblock Vol.~28, Monographs in Population Biology, Princeton University
  Press, New Jersey.

\bibitem[ICES(2015)]{ices15}
ICES (2015).
\newblock {Manual for the International Bottom Trawl Surveys}.
\newblock Series of ICES Survey Protocols SISP 10 -- IBTS IX. 86~p.

\bibitem[Jennings(2005)]{jenn05}
Jennings, S. (2005).
\newblock Indicators to support an ecosystem approach to fisheries.
\newblock {\em Fish Fish.}, {\bf 6}, 212--232.

\bibitem[Jennings and Dulvy(2005)]{jd05}
Jennings, S. and Dulvy, N.~K. (2005).
\newblock Reference points and reference directions for size-based indicators
  of community structure.
\newblock {\em ICES J.~Mar.~Sci.}, {\bf 62}, 397--404.

\bibitem[Kelly-Gerreyn {\em et~al.}(2014)]{kellyg14}
Kelly-Gerreyn, B.~A., Martin, A.~P., Bett, B.~J., Anderson, T.~R., Kaariainen,
  J.~I., Main, C.~E., Marcinko, C.~J., and Yool, A. (2014).
\newblock Benthic biomass size spectra in shelf and deep-sea sediments.
\newblock {\em Biogeosciences}, {\bf 11}, 6401--6416.

\bibitem[Lawless(2003)]{lawl03}
Lawless, J.~F. (2003).
\newblock {\em Statistical Models and Methods for Lifetime Data}.
\newblock 2nd ed., Wiley series in Probability and Statistics, Wiley, New
  Jersey.

\bibitem[Love {\em et~al.}(2002)]{lyt02}
Love, M.~S., Yoklavich, M., and Thorsteinson, L. (2002).
\newblock {\em The Rockfishes of the Northeast Pacific}.
\newblock University of California Press, Berkeley and Los Angeles, California.

\bibitem[Maxwell and Jennings(2006)]{mj06}
Maxwell, T. A.~D. and Jennings, S. (2006).
\newblock Predicting abundance-body size relationships in functional and
  taxonomic subsets of food webs.
\newblock {\em Oecologia}, {\bf 150}, 282--290.

\bibitem[McCoy {\em et~al.}(2016)]{mcco16}
McCoy, K., Heenan, A., Asher, J., Ayotte, P., Gorospe, K., Gray, A., Lino, K.,
  Zamzow, J., and Williams, I. (2016).
\newblock {Pacific Reef Assessment and Monitoring Program data report.
  Ecological monitoring 2015 -- reef fishes and benthic habitats of the main
  Hawaiian Islands, Northwest Hawaiian Islands, Pacific Remote Island Areas,
  and American Samoa}.
\newblock PIFSC Data Report DR-16-002.

\bibitem[Mindel {\em et~al.}(2018)]{mnwb18}
Mindel, B.~L., Neat, F.~C., Webb, T.~J., and Blanchard, J.~L. (2018).
\newblock Size-based indicators show depth-dependent change over time in the
  deep sea.
\newblock {\em ICES J.~Mar.~Sci.}, {\bf 75}, 113--121.

\bibitem[Monnahan {\em et~al.}(2016)]{monn16}
Monnahan, C.~C., Ono, K., Anderson, S.~C., Rudd, M.~B., Hicks, A.~C.,
  Hurtado-Ferro, F., Johnson, K.~F., Kuriyama, P.~T., Licandeo, R.~R., Stawitz,
  C.~C., Taylor, I.~G., and Valero, J.~L. (2016).
\newblock The effect of length bin width on growth estimation in integrated
  age-structured stock assessments.
\newblock {\em Fish.~Res.}, {\bf 180}, 103--112.

\bibitem[Piet and Jennings(2005)]{pj05}
Piet, G.~J. and Jennings, S. (2005).
\newblock Response of potential fish community indicators to fishing.
\newblock {\em ICES J.~Mar.~Sci.}, {\bf 62}, 214--225.

\bibitem[Platt and Denman(1978)]{pd78}
Platt, T. and Denman, K. (1978).
\newblock The structure of pelagic marine ecosystems.
\newblock {\em Rapp. P.-v. R{{\'e}}un. Cons. Int. Explor. Mer}, {\bf 173},
  60--65.

\bibitem[Rice and Gislason(1996)]{rg96}
Rice, J. and Gislason, H. (1996).
\newblock Patterns of change in the size spectra of numbers and diversity of
  the {N}orth {S}ea fish assemblage, as reflected in surveys and models.
\newblock {\em ICES J.~Mar.~Sci.}, {\bf 53}, 1214--1225.

\bibitem[Rice {\em et~al.}(2012)]{rice12}
Rice, J., Arvanitidis, C., Borja, A., Frid, C., Hiddink, J.~G., Krause, J.,
  Lorance, P., Ragnarsson, S.~A., Sk\"old, M., Trabucco, B., Enserink, L., and
  Norkko, A. (2012).
\newblock {Indicators for sea-floor integrity under the European Marine
  Strategy Framework Directive}.
\newblock {\em Ecol.~Indic.}, {\bf 12}, 174--184.

\bibitem[Richards {\em et~al.}(2011)]{rwnz11}
Richards, B.~L., Williams, I.~D., Nadon, M.~O., and Zgliczynski, B.~J. (2011).
\newblock A towed-diver survey method for mesoscale fishery-independent
  assessment of large-bodied reef fishes.
\newblock {\em Bull.~Mar.~Sci.}, {\bf 87}, 55--74.

\bibitem[Robinson {\em et~al.}(2017)]{robins17}
Robinson, J. P.~W., Williams, I.~D., Edwards, A.~M., McPherson, J., Yeager, L.,
  Vigliola, L., Brainard, R.~E., and Baum, J.~K. (2017).
\newblock Fishing degrades size structure of coral reef fish communities.
\newblock {\em Glob.~Change~Biol.}, {\bf 23}, 1009--1022.

\bibitem[Robinson {\em et~al.}(2010)]{robins10}
Robinson, L.~A., Greenstreet, S. P.~R., Reiss, H., Callaway, R., Craeymeersch,
  J., {de Boois}, I., Degraer, S., Ehrich, S., Fraser, H.~M., Goffin, A.,
  Kr\"oncke, I., {Lindal Jorgenson}, L., Robertson, M.~R., and Lancaster, J.
  (2010).
\newblock {Length-weight relationships of 216 North Sea benthic invertebrates
  and fish}.
\newblock {\em J. Mar. Biol. Assoc. U.K.}, {\bf 90}, 95--104.

\bibitem[Scott {\em et~al.}(2014)]{sba14}
Scott, F., Blanchard, J.~L., and Andersen, K.~H. (2014).
\newblock \emph{mizer}: an {R} package for multispecies, trait-based and
  community size spectrum ecological modelling.
\newblock {\em Meth.~Ecol.~Evol.}, {\bf 5}, 1121--1125.

\bibitem[Sheldon and Kerr(1972)]{sk72}
Sheldon, R.~W. and Kerr, S.~R. (1972).
\newblock {The population density of monsters in Loch Ness}.
\newblock {\em Limnol.~Oceanogr.}, {\bf 17}, 796--798.

\bibitem[Shin {\em et~al.}(2005)]{srjfg05}
Shin, Y.-J., Rochet, M.-J., Jennings, S., Field, J.~G., and Gislason, H.
  (2005).
\newblock Using size-based indicators to evaluate the ecosystem effects of
  fishing.
\newblock {\em ICES J.~Mar.~Sci.}, {\bf 62}, 384--396.

\bibitem[Stasko {\em et~al.}(2016)]{ssmarp16}
Stasko, A.~D., Swanson, H., Majewski, A., Atchison, S., Reist, J., and Power,
  M. (2016).
\newblock {Influences of depth and pelagic subsidies on the size-based trophic
  structure of Beaufort Sea fish communities}.
\newblock {\em Mar.~Ecol.~Prog.~Ser.}, {\bf 549}, 153--166.

\bibitem[Taniguchi {\em et~al.}(2014)]{tfp14}
Taniguchi, D. A.~A., Franks, P. J.~S., and Poulin, F.~J. (2014).
\newblock Planktonic biomass size spectra: an emergent property of
  size-dependent physiological rates, food web dynamics, and nutrient regimes.
\newblock {\em Mar.~Ecol.~Prog.~Ser.}, {\bf 514}, 13--33.

\bibitem[Thorpe {\em et~al.}(2015)]{tllcj15}
Thorpe, R.~B., {Le~Quesne}, W. J.~F., Luxford, F., Collie, J.~S., and Jennings,
  S. (2015).
\newblock Evaluation and management implications of uncertainty in a
  multispecies size-structured model of population and community responses to
  fishing.
\newblock {\em Meth.~Ecol.~Evol.}, {\bf 6}, 49--58.

\bibitem[Trebilco {\em et~al.}(2015)]{tdss15}
Trebilco, R., Dulvy, N.~K., Stewart, H., and Salomon, A.~K. (2015).
\newblock The role of habitat complexity in shaping the size structure of a
  temperate reef fish community.
\newblock {\em Mar.~Ecol.~Prog.~Ser.}, {\bf 532}, 197--211.

\bibitem[{van~Gemert} and Andersen(2018)]{va18}
{van~Gemert}, R. and Andersen, K.~H. ({2018}).
\newblock Implications of late-in-life density-dependent growth for fishery
  size-at-entry leading to maximum sustainable yield.
\newblock {\em ICES J.~Mar.~Sci.}, {\bf 75}, 1296--1305.

\bibitem[White {\em et~al.}(2007)]{weke07}
White, E.~P., Ernest, S. K.~M., Kerkhoff, A.~J., and Enquist, B.~J. (2007).
\newblock Relationships between body size and abundance in ecology.
\newblock {\em Trends Ecol.~Evol.}, {\bf 22}, 323--330.

\bibitem[White {\em et~al.}(2008)]{weg08}
White, E.~P., Enquist, B.~J., and Green, J.~L. (2008).
\newblock On estimating the exponent of power-law frequency distributions.
\newblock {\em Ecology}, {\bf 89}, 905--912.

\bibitem[Woodson {\em et~al.}(2018)]{wsj18}
Woodson, C.~B., Schramski, J.~R., and Joye, S.~B. (2018).
\newblock A unifying theory for top-heavy ecosystem structure in the ocean.
\newblock {\em Nat. Commun.}, {\bf 9}, 23.

\bibitem[Yool {\em et~al.}(2017)]{ymabjr17}
Yool, A., Martin, A.~P., Anderson, T.~R., Bett, B.~J., Jones, D. O.~B., and
  Ruhl, H.~A. (2017).
\newblock Big in the benthos: Future change of seafloor community biomass in a
  global, body size-resolved model.
\newblock {\em Glob.~Change~Biol.}, {\bf 23}, 3554--3566.

\bibitem[Zhang {\em et~al.}(2018)]{zcxxr18}
Zhang, C., Chen, Y., Xu, B., Xue, Y., and Ren, Y. (2018).
\newblock Evaluating fishing effects on the stability of fish communities using
  a size-spectrum model.
\newblock {\em Fish.~Res.}, {\bf 197}, 123--130.

\bibitem[Zhou {\em et~al.}(2019)]{zhou19}
Zhou, S., Kolding, J., Garcia, S.~M., Plank, M.~J., Bundy, A., Charles, A.,
  Hansen, C., Heino, M., Howell, D., Jacobsen, N.~S., Reid, D.~G., Rice, J.~C.,
  and {van~Zwieten}, P. A.~M. (2019).
\newblock Balanced harvest: concept, policies, evidence, and management
  implications.
\newblock {\em Rev.~Fish Biol.~Fish.}, {\bf 29}, 711--733.

\end{thebibliography}


\clearpage

%%% TABLE 1 %%%
% Species names
% latex table generated in R 3.4.3 by xtable 1.8-2 package
% Tue May 15 15:12:05 2018
\begin{table}[tp]
\centering
\caption{First six and last six rows of the IBTS data to illustrate the
  information available. Each row represents a unique combination of year,
  species and length class. `Number' gives the number of individuals per hour of
  trawling observed for that combination, and can be non-integer because counts
  of individual fish are scaled by tow duration. Fish lengths are assigned into
  length classes, where the value of `Length class' (cm) is the minimum value of
  the 1-cm length bin (or 0.5-cm length bins for Atlantic Herring and European
  Sprat). Parameters $\alpha_s$ and $\beta_s$ are the length-weight coefficients
  for species $s$ from \citet{ffrr12}, and `Body mass' (g) is the resulting
  estimated body mass for an individual of that species and length
  class. `Biomass' (g~h$^{-1}$) is the total biomass caught per hour of trawling
  for each row. Example species are Smallspotted Catshark ({\it Scyliorhinus
    canicula}), Snakeblenny ({\it Lumpenus lampretaeformis}) and Thickback Sole
  ({\it Microchirus variegatus}). The full data set has 42,298 rows.}
\label{tab:egData}
\begin{tabular}{rrrrrrrr}
  \hline
  Year & Species & Length & Number & $\alpha_s$ & $\beta_s$ & Body & Biomass\\
       & & class (cm) & (h$^{-1}$) & & & mass (g) & (g~h$^{-1}$) \\
  \hline
  1986 & Smallspotted Catshark & 45 & 0.007 & 0.0031 & 3.0290 & 315.46 & 2.25 \\
  1986 & Smallspotted Catshark & 46 & 0.007 & 0.0031 & 3.0290 & 337.17 & 2.41 \\
  1986 & Smallspotted Catshark & 50 & 0.007 & 0.0031 & 3.0290 & 434.05 & 3.10 \\
  1986 & Smallspotted Catshark & 52 & 0.029 & 0.0031 & 3.0290 & 488.81 & 14.33 \\
  1986 & Smallspotted Catshark & 53 & 0.011 & 0.0031 & 3.0290 & 517.84 & 5.65 \\
  1986 & Smallspotted Catshark & 54 & 0.011 & 0.0031 & 3.0290 & 548.00 & 6.18 \\
  $\cdot \cdot \cdot$ & $\cdot \cdot \cdot$ & $\cdot
  \cdot \cdot$ & $\cdot \cdot \cdot$ & $\cdot \cdot \cdot$ & $\cdot \cdot \cdot$
               & $\cdot \cdot \cdot$ & $\cdot \cdot \cdot$ \\
  2015 & Snakeblenny & 34 & 0.028 & 0.0244 & 2.0439 & 32.93 & 0.92 \\
  2015 & Thickback Sole & 8 & 0.013 & 0.0080 & 3.1410 & 5.49 & 0.07 \\
  2015 & Thickback Sole & 14 & 0.039 & 0.0080 & 3.1410 & 31.85 & 1.24 \\
  2015 & Thickback Sole & 15 & 0.052 & 0.0080 & 3.1410 & 39.55 & 2.05 \\
  2015 & Thickback Sole & 16 & 0.065 & 0.0080 & 3.1410 & 48.44 & 3.15 \\
  2015 & Thickback Sole & 17 & 0.013 & 0.0080 & 3.1410 & 58.60 & 0.76 \\
  \hline
\end{tabular}
\end{table}

%%% TABLE 2 %%%
\begin{table}[tp]
\centering
\caption{Recommendations for calculating the size-spectrum exponent,
  $b$, of the individual size distribution for body masses (equation \ref{PLB}),
  depending on the type and resolution of the data.}
\label{tab:recommend}
\small{
\begin{tabular}{L{3 cm}L{3 cm}L{5 cm}L{4 cm}}
\hline
Data type & Resolution & Issue & Solution \\
\hline
  Body mass & Individual body masses & Accuracy of measurements may require them
                                       to be considered as bins
          & Simulate data as in Figure~\fitMLEhists~to see if MLE method
            sufficient or MLEbin is needed\\
Body mass & Individuals assigned to common body-mass bins & Data are counts in bins
                               & Use MLEbin method\\
Length & Individual lengths & Lengths have to be converted to body masses
       & Calculate individual body masses and proceed as per `individual body masses'\\
Length & Individuals assigned to length bins &
         Resulting converted body-mass bins are species specific &
         Use MLEbins method\\
Lengths from underwater visual census (e.g.~point counts or belt transects)
       & Individual lengths
       & May need to standardise counts when combining different survey designs
         (e.g.~\citealt{heen16}); surveys may cover different sized areas or
         count small and large organisms differently
       & Potential need to adapt MLEbins likelihood function to account for particular
         survey design \\
\hline
\end{tabular}
}
\end{table}

\renewcommand{\baselinestretch}{1.1}

\onefigshift{nSeaFung/nSeaFungTrends}{For the IBTS data, each of eight methods is
  used to estimate the exponent $b$ (circles, with 95\% confidence
  interval as vertical bars) for each year. The fit of a weighted linear
  regression with 95\% confidence interval is shown in red if the trend
  can be considered statistically significant from 0 ($p<0.05$) and in grey if
  it cannot be ($p \geq 0.05$).
  The y-axes are the same for (d)-(h). For (g) and (h) the confidence
  intervals are too narrow to be seen. See Supplementary Material for full
  details of methods.}{figure1}


\onefigshift{binFigs/binLW}{Example length-weight relationships (solid curves)
  for Lemon Sole ({\it Microstomus kitt}) in blue and Common Ling ({\it Molva
    molva}) in red, showing how binned length measurements yield ranges of
  estimated body masses. Horizontal bars show six example 5-cm length bins, the
  same bins for both species. Using the length-weight relationships, these length
  bins convert to the body-mass bins shown by the vertical bars, giving six bins
  of unequal width for each species. The conversion for the 30-35~cm length bin
  is explicitly shown as dotted lines for the endpoints and solid lines for the
  midpoint of 32.5~cm. For Lemon Sole, the 30-35~cm bin with midpoint 32.5~cm
  converts to a body-mass bin of 309-473~g (values
  rounded to nearest~g for clarity), and for Common Ling the converted body-mass
  bin is 119-202~g.
  Note that converting the midpoint of 32.5~cm gives respective body masses of
  385~g and 157~g, which are slightly
  less than the respective midpoints of the new body-mass bins (391~g and
  161~g) because of the nonlinearity of the length-weight relationships.}{figure2}

\clearpage

\onefig{binFigs/binLWlog}{Demonstration of how binned body-mass values
  (blue bars) get
  assigned to logarithmic size-class bins (black and grey bars),
  as occurs for the LBbiom and LBNbiom
  methods that prescribe the logarithmic bins. Blue vertical bars show the
  resulting six converted body-mass bins for Lemon Sole from
  Figure~\ref{fig:binFigs/binLW}. The black and grey vertical bars show
  body-mass bins that have equal width on a logarithmic scale. Explicitly these
  are at $4, 8, 16, 32, ...$, as used in the LBbiom and LBNbiom methods. The
  short horizontal purple lines show the midpoints of those bins. Each panel,
  labelled Bin~1 to Bin~6, highlights one blue body-mass bin. For each blue bin,
  the blue arrow starts from the converted value of the midpoint of the original
  length bin (see Figure~\ref{fig:binFigs/binLW} for Bin~5), and ends at the
  purple midpoint of the $\logtwo$ bin in which the converted value lies. The
  blue dashed lines indicate that all body masses within the blue bin get
  assigned to the same purple midpoint value. This ignores uncertainties in
  body masses, and for Bins~3 and~6 assigns completely incorrect values
  (greater than the body mass of any individual within the bin). But
  such values get
  used in the LBbiom or LBNbiom methods.}{figure3}

\clearpage

\thispagestyle{empty}

\onefig{binFitting/fitMLEhists}{Histograms of estimated exponent $b$ for 10,000
  simulated data sets, each of which contains 1,000 independent random numbers
  drawn from a bounded power-law distribution with $b=-2$, $\xmin=1$ and
  $\xmax=1,000$.
Panels (a)-(f) use linear bins of constant bin width 1, 5 or 10 (binning type
`Linear 1' etc.), and panels (g)-(h) use bin widths that double in size (binning
type `2k'). Each binned data set is fitted using either the MLEmid method (a, c,
e, g) or the MLEbin method (b, d, f, h). Vertical red lines indicate the known
value of $b=-2$.}{figure4}

\clearpage
\thispagestyle{empty}

\onefigshift{binFitting/fitMLEconfs}{Confidence intervals (horizontal lines) of
  $b$ obtained for each combination of data-binning and MLE method for
  subsamples of the 10,000 simulated data sets (with $\xmax=1,000$) used in
  Figure~\fitMLEhists. Panels are the same as in Figure~\fitMLEhists. For each
  numbered subsample on the y-axis, the 95\% confidence interval of $b$ obtained
  using the respective method is plotted as a horizontal line, which is coloured
  grey if the interval includes the true value of $b=-2$ (given by the vertical
  red line) or blue if it does not. Simulations are sorted in ascending order of
  their lower bound. The percentage for each method gives the observed coverage,
  namely the percentage of all 10,000 simulated data sets for which the 95\%
  confidence interval contains the true value of $b$; by definition, this should
  ideally be 95\%. Horizontal axes are the same for all panels here and in
  Figure~\fitMLEhists.}{figure5}

\clearpage

\onefig{nSeaFung/nSeaFungBins1}{Vertical bars show the body-mass bins that
  result from converting the 1-cm length bins into body-mass bins using the
  species-specific length-weight parameters. These are the first 45 species of
  the IBTS data, ordered by the maximum value of the highest bin. Atlantic
  Herring and European Sprat use 0.5-cm length bins (indicated by~$\times$). Any
  vertical gaps between vertical bars indicate that there were no body masses of
  that species in that range; i.e.~there were no individuals in the length
  bin(s) that would result in the body-mass bin(s) that would fill the gap. For
  example, there were no $\sim$30~g European Sprat (code 126425). Two example species are
  shown in red. They have similar ranges of body masses, but Moustache Sculpin
  (code 127205, \emph{Triglops murrayi}) is short and fat with lengths from
  8-15~cm and length-weight relationship $w = 0.0088~l^{3.00}$, whereas
  Snakeblenny (code 154675, \emph{Lumpenus lampretaeformis}) is long and skinny
  with lengths from 13-37~cm and $w = 0.0244~l^{2.04}$, and consequently the
  1-cm length bins translate to much narrower body-mass bins than for Moustache
  Sculpin. Equivalent figures for the remaining 90 species are given in the
  Supplementary Material.}{figure6}


\clearpage
\thispagestyle{empty}
\begin{figure}[tp]
\begin{center}
\vspace{-14mm}
\epsfbox{figure7.eps} % ../../fitting-size-spectra/codeBinStructure/nSeaFung/IBTS-ISDfigs/IBTS-ISD1999.eps}
\end{center}
\vspace{-9mm}
\caption{Suggested plots of individual size distribution and MLEbins fit
            (red solid curve)
            with 95\% confidence intervals (red dashed curves, may be hard to
            see) for IBTS data in 1999. The y-axis is linear in (a) and logarithmic
            in (b).
            For each bin, the horizontal green line shows
            the range of body sizes,
            with its value on the y-axis corresponding to the
            total number of individuals (per hour of trawling) in bins whose
            minima are $\geq$ the
            bin's minimum.
            The vertical span of each grey rectangle shows the
            possible range of the number of individuals with body mass
            $\geq$ the body mass of individuals in that bin (horizontal span
            is the same as for the green lines).
            The text in (a) gives the year,
            the MLE for the size-spectrum exponent $b$, and the sample size $n$.
            See Supplementary Material for
            full details and results for all years.
            Note that in (b) the logarithmic y-axis compresses most of the data
            into the top-left corner.}
\label{fig:ISD-1999.main}
\end{figure}
\clearpage

\begin{figure}[tp]
\begin{center}
\epsfxsize=6in
\epsfbox{figure8.eps}
  \end{center}
  \caption{Comparison of the annual estimates of
  $b$ for the IBTS data using the original MLE method and the MLEbins method,
  showing medians and 95\% confidence intervals (each year is slightly offset
  to clearly show the confidence intervals).}
\label{fig:nSeaFung/nSeaFungCompareTrends}
\end{figure}

\clearpage

\end{document}
